\documentclass{article}
\title{Two's Complement Encodings}
\author{CSAPP}
\usepackage{booktabs}
\usepackage{physics}
\usepackage{amsmath}
\usepackage{parskip}
\begin{document}
\maketitle
\begin{flushleft}
\textbf{Definition}\\
    Definition of two's complement encoding.\\
    For vector
    \[\vec{x} = [x_{w-1}, x_{w-2}, \ldots, x_{0}]\]
    \[B2T_w(\vec{x}) \doteq -x_{w-1} \cdot 2^{w-1} + \sum_{i=0}^{w-2} x_{i} \cdot 2^{i}\]

\textbf{Note}\\
    The most significant bit $x_{w-1}$ is also called the sign bit. Its
    weight is $-2^{w-1}$, the negation of its weight in an unsigned representation.
    When the sign bit is set to 1, the represented value is negative, and when set to
    0, the value is nonnegative.

\smallskip
\textbf{Range}\\
    Two's complement representation has an asymmetric range. For a word size of $w$ bits,
    the minimum and maximum representable values are defined as:
    \[TMin_w \doteq -2^{w-1}\]
    \[TMax_w \doteq 2^{w-1} - 1\]
    For example, with $w = 4$:
    \[TMin_4 = -2^3 = -8\]
    \[TMax_4 = 2^3 - 1 = 7\]

\textbf{Note}\\
    Observe that $|TMin_w| = TMax_w + 1$. That is, the range of negative numbers is
    one larger than the range of positive numbers. This asymmetry arises because zero
    is considered nonnegative, so the positive range loses one value. As a consequence,
    arithmetic operations on two's complement numbers can sometimes lead to overflow
    at the negative end without a symmetric positive counterpart. Note also that
    $TMax_w + 1 = TMin_w$ under two's complement arithmetic, which is a common
    source of subtle bugs in systems programming.

\smallskip
\textbf{Sample}\\
    \[B2T_{4}(0001) = -0\cdot2^3 + 0\cdot2^2 + 0\cdot2^1 + 1\cdot2^0 = 0+0+0+1 = 1\]
    \[B2T_{4}(0101) = -0\cdot2^3 + 1\cdot2^2 + 0\cdot2^1 + 1\cdot2^0 = 0+4+0+1 = 5\]
    \[B2T_{4}(1011) = -1\cdot2^3 + 0\cdot2^2 + 1\cdot2^1 + 1\cdot2^0 = -8+0+2+1 = -5\]
    \[B2T_{4}(1111) = -1\cdot2^3 + 1\cdot2^2 + 1\cdot2^1 + 1\cdot2^0 = -8+4+2+1 = -1\]
\end{flushleft}
\end{document}